\documentclass[11pt]{article}
\usepackage[a4paper,margin=22mm]{geometry}
\usepackage{fontspec}
\setmainfont{LinLibertine_R.otf}[BoldFont=LinLibertine_RB.otf,ItalicFont=LinLibertine_RI.otf,BoldItalicFont=LinLibertine_RBI.otf]
\usepackage{xcolor}
\usepackage{hyperref}
\hypersetup{colorlinks=true,urlcolor=blue}
\setlength{\parindent}{0pt}
\setlength{\parskip}{6pt}
\sloppy
\emergencystretch=2em
\begin{document}
\section*{Historical Britannica: extracted source}
Encyclopaedia Britannica - Third edition, Volume 5, CIC-DIA - EB.5\par
{\small Source row 1230. Parser: nvidia/nemotron-3.5-parse.\par}
{\small This is model-extracted text, not a corrected transcription. OCR errors, reading order and region alignment require human review.\par}
\hrule\medskip
{\small\itshape C L A [ 47 ] C L A\par}

{\scriptsize\color{gray} Region region-0; provider order 0.\par}
\medskip

{\small\itshape **Clay.**\par}

{\scriptsize\color{gray} Region region-1; provider order 1.\par}
\medskip

M. Beume alfo pretends that clay contains a little vitrile acid, and is therefore fobble in a large quantity of boiling water. Mr Schebel likewise tried this experiments but found, that of several kinds of argyliaceous earth not the finalelfi quantity was diffused; and an hu fewife made feared experiments in order obtain virkleoid. from pure clay, line without fluo- cefs: neither was he able to obtain any lepar either by means of all other of tartar or with removal) nor could be obtain with clay a strindineutral fat from the refriduum of the dithilation of muricans and ferrous acid.
\par

{\scriptsize\color{gray} Region region-2; provider order 2.\par}
\medskip

The correction of clay by heat has affected Mr Wedgewood an opportunity of conflicting by Jia necan in information for mealtering those degrees of heat which are above the most of the fate of common thermometers, as deferred under the article TrasvSYvers.
\par

{\scriptsize\color{gray} Region region-3; provider order 3.\par}
\medskip

Mr Schebel has made feered experiments to differver the properties of alum when mixed with other fullerenes. A Solution of alum, he finds, is decon- points by limonwater, and if no more of the water be added than Is exactly requisite for the precipitation, the earth of alum forms a treatment precipitate like boiled freshly, and if the shear water be filtered, it is found to be a foliation of gypfum. On adding more lime-water than a secretory for precipitation, the pro-ipitated barrier is almost of the gelatinates appara- nace jult mentioned. If the whole be allowed to fland for a quarter of an hour, and frequently adjusted during that time, too pygmy area, even lines, is found in the filtered liquor, unless too much lime-water has been used. On examining the precipitate in this effect, our author found it to could end of some, fetonic, and lime. This was discovered frill by treating it with miruistic acid, which differed the aluminum earth, leaving the gapion behind. The addition of confine volatile alkali throw down a transparent gelatinations may, which on the earth of alum, i and on finishing it again, and then adding a fixed alkaline lisivium, the lime was thrown down y wherece it appeared, that the lime and gyprum had separated from the water, and united with the earth of alum.
\par

{\scriptsize\color{gray} Region region-4; provider order 4.\par}
\medskip

To understand the reason of this uncommon preci- ptition. Mr Schebe next poured into a foliation of alum a quantity of earlier volatile alkali more than fufect or fauture the acid, in order to be certain having it all taken at. The precipitated redan way then edolcorated, and mixed with a foliation of gypfum, that he might coherent whether the gypsum would fo- peate from the water, and precipitate with the earth of alum; which, however, did not take place. On moving lime water with the peppitas, he found that the former very dense hull in smaller sfe, and that the earth of alum become optique. Some part of the wa-ter was drained, ilavimu tarter dropped into it; but it remained clear, now one any proportion formed by a foliation of corrosive fuhlinate. He afterwashed metallic soil to the fall peppitas, which it differed entirely without leaving any gygrum be- kind; whence our author concludes, that the earth of alum halted invited with the lime into a peculiar kind of compound.
\par

{\scriptsize\color{gray} Region region-5; provider order 5.\par}
\medskip

Lastly, he now imagined, that this compound of earth of alum and lime might be capable of separating gypfum from water. To try this, be prepared a large
\par

{\scriptsize\color{gray} Region region-6; provider order 6.\par}
\medskip

quantity of the compound earth, mixed it with a fola- tion of gypfum, and let it refl for a quarter of an hour when he found, so the farphit, that the gypforn fill imunited infepedel in the water, and that the pre- pitate was entirely fobble in muriatic acid. He now rinses to foliated of gypfum with lime-water, adding earth of alum at the home times), when he found, that the whole we presticipated as before, the lime and ala- minicous earth having fallen on the booking with the gypfum, leaving the water gone. On the whole, our author concludes, i That the vitrile acid in gyprus is capable of combining with more lime than as secretary to an exact futuration. 2. The calcareous earth capable of forming an unison with the earth of alum, i But theygripe cannot combine with the earth of alum; but that if a paperflaun quality of lime the united with virleid acid, it will then there as blood of unam a combine grydim with the earth of alum, and thus form a new compound confiling of these cartels. Pure day has so effect upon limonwater.
\par

{\scriptsize\color{gray} Region region-7; provider order 7.\par}
\medskip

Concified is of opinion, that common clay especially the blue, grey, and red kinds, may derive their origin from mud (and as the man proceeded from vegetables, will have fallen, that the varieties of cup) mentioned are nothing else but the common mould al- terde, after a length of time, by means o water. The option, he thanks, receives considerable through from the following circumferences : vie, that a great quantity of fireplants not every year in the lakes, and changed into mud syry hold of which, however, is been upon the floorer after the water is dried in the former times; that that the degree where the mud works. Professor Bergen has likewise indeed, that pure clay may be a calcereous earth combined with home and not yet differences \textbackslash{}* We (days he) compositions of this kind ought to be considered in premige fila- distances, with respect to our knowledge of them, till my fault by experimentally decomposed, for no found knowledge in natural philosophy can be obtained from the consideration of mere polibilities / force daily ex- perverse flows, that even the such probable fogpolloious law proved falls, when the means of putting them to the tell how afterwards been found out."— Now, therefore, (days M. Mingallah), that the \_argilliaceae\_ is acknowledged to be a simple primitive earth, which cannot be decoupled into any other principles, the format by the combination of any other simple fulhifances we know, we ought to set, finished at, pre- tentwithout endeavouring to account for itsformation.”
\par

{\scriptsize\color{gray} Region region-8; provider order 8.\par}
\medskip

The principal figures of the argillaceous earth or clays are.
\par

{\scriptsize\color{gray} Region region-9; provider order 9.\par}
\medskip

3. The \_avgila arron\_, or the base. It is generally found in much cakes of the differing of shale, like which, alto, it marks white. Its harducls is nearly like that of the forests, and it faced his first than days com- only. No. In of all sixteenth edition, and about the specific gravity of 1.66g. When examined with a mi- envelope, it is found a cotif of full transparent cry- thicks and, from Mr Scheber’s aspirations, appears to be an argillaceous earth flattered with fixed air, in consequence of which it otherwise with acids. It contains also a finally up to calcerements, and femmes of gypfum, with fome flight traces of iron. It is found at Italales.
\par

{\scriptsize\color{gray} Region region-10; provider order 10.\par}
\medskip

\#\# 2. The orgula agres, porcelain clay, the kaolin of the
\par

{\scriptsize\color{gray} Region region-11; provider order 11.\par}
\medskip
\newpage\section*{Source record}
bri-1230-4f96d9bebe42\par
Image SHA-256: \texttt{4f96d9bebe421b6d59e2a525f5b1a686}\par\texttt{601a8468f7bf92ed458f450ebd6e14b3}\par
Dataset revision: \texttt{309f765b1ecc83e0a7f7}\texttt{762e44b4a927a7555443}\par
National Library of Scotland, Encyclopaedia Britannica Illustrated (CC0).\par
\url{https://huggingface.co/datasets/NationalLibraryOfScotland/encyclopaedia_britannica_illustrated/viewer/default/train?row=1230}
\end{document}
